\section{Methods}
\label{sec:methods}

\subsection{Tasks and calibration scenarios}
\label{sec:scenarios}
We studied three tasks: rest against hand against sixth finger (three
classes; chance is one third), sixth finger against rest (detecting the new
command) and hand against sixth finger (the published task). Calibration was
simulated on the recorded sessions. For each target patient we drew $k$
trials of each imagery class, with their rest windows, from runs 1--2, and
tested on runs 3--4 of the same session, which no step ever saw. The budgets
were $k \in \{0, 1, 2, 3, 5, 10, \text{all}\}$, with ten random draws per
budget (five for the networks) shared by every method. Minutes are given at
the article's nominal 14\,s per trial: one trial of each movement costs
\rqTwoMinutesKOne\,min and the whole pool \rqTwoMinutesKAll\,min. At the
recorded trial length they are \rqTwoMinutesRecordedKOne{} and
\rqTwoMinutesRecordedKAll\,min.

In the \emph{new-user} scenario the patient records $k$ hand and $k$
sixth-finger trials. In the \emph{new-command} scenario the patient's hand
trials of runs 1--2 are already recorded and $k$ sixth-finger trials are
added. Its three-class decoders learn from every hand trial and the $k$
sixth-finger trials, with their rest windows. Its sixth-finger-against-rest
decoders, as pre-registered, learn only from the sixth-finger and rest
windows of the $k$ new trials and use no hand trial; for that task the two
scenarios differ only in their random draws, and neither adds a command to a
hand decoder.

\subsection{Decoders}
The confirmatory pair, fixed before any result, compares one classifier with
and without other patients. \emph{From scratch} (TS+LR) estimates a covariance
matrix per window, projects it to the tangent space (TS) of the Riemannian
manifold and fits a logistic regression (LR) to the patient's
windows~\cite{barachant2012}. With one trial of each movement it has four
training windows, two of them rest. In the new-command scenario its classes
are unbalanced by design (up to 20 hand trials against $k$ sixth-finger
trials), and its logistic regression, as pre-registered, has no class
weights.

\emph{Transfer} (RCT-TS+LR) first re-centres each patient's covariance
matrices on that patient's own Riemannian mean, so that all patients share
one reference point~\cite{zanini2018}. It then fits the same classifier to the
other patients' windows together with the target's. The target's mean comes
from its drawn windows; at $k = 0$ it comes from the rest windows of runs
1--2 (\rqTwoRestMinutes\,min of unlabelled recording, nominal). The weight of
the target's windows was chosen once by the pre-registered rule, on two
development patients without touching any runs 3--4, from \{1, 5, 20\}. All
three candidates scored near chance (\rqTwoHparamRctTslrMeanMin--\rqTwoHparamRctTslrMeanMax\,\%
on three classes), and the highest-scoring one was the smallest,
\rqTwoHparamRctTslrWeight. Each of the target's at most 80 windows therefore
counts as much as one window of the other patients, who contribute thousands
(all their sessions; the dataset has \nWindows{} windows).

The secondary pair compares EEGNet~\cite{lawhern2018} trained on the patient
alone with an EEGNet pre-trained on the other patients after Euclidean
alignment~\cite{he2020} and then fine-tuned on the patient. The same rule
chose its fine-tuning values from a learning rate of $10^{-4}$ or
$3\times10^{-4}$ and 20 or 50 epochs; the highest-scoring pair was again the
least adaptive, a rate of \rqTwoHparamEegnetLr{} and \rqTwoHparamEegnetEpochs{}
epochs. For the new
command, a control called the \emph{own-hand decoder} is TS+LR trained on the
patient's hand and rest trials only, with any non-rest output counted as the
sixth finger. CSP with linear discriminant analysis
(CSP+LDA)~\cite{ramoser2000}, minimum distance to the Riemannian mean
(MDM)~\cite{barachant2012}, TS+LR on pooled patients without re-centring, and
re-centred MDM ran as descriptive references. Source data were every window
of every other patient, never the target's linked partner
(Section~\ref{sec:dataset}). The selection of the transfer values used the
same sources, so other patients' runs 3--4 took part in choosing values that
were later tested on those runs, a weak breach of nesting; the candidates
scored within a point of each other.

\subsection{Leakage-free evaluation}
The dataset's classification script fits the CSP filters on all trials of a
session and only then runs 10-fold cross-validation of the
classifier~\cite{wang2026dataset}. We rebuilt this protocol and compared it
with the same 10$\times$10 folds in which every step is fitted inside the
training folds (leakage-free cross-validation), and with two run-wise splits,
leave one run out and runs 1--2 against 3--4 (chronological). Hypothesis H1
tests the inflation for CSP+LDA on hand against sixth finger. H1b repeats the
test with the same code on two public healthy-subject datasets,
BNCI2014\_001~\cite{tangermann2012} and Cho2017~\cite{cho2017}, loaded with
MOABB~\cite{jayaram2018moabb}. All other results use the chronological split.
Before every fit the code checked that no test trial, and no window of the
target or its linked partner, entered the training or source data, and
stopped otherwise.

\subsection{Metrics and statistics}
The metric is balanced accuracy, with the patient as the unit; differences in
accuracy are in percentage points. Each patient's calibration curve is
summarised by a normalised area under the learning curve (ALC), after the
active learning challenge~\cite{guyon2011}:
\begin{equation}
\mathrm{ALC} = \frac{1}{x_{\max} - x_{\min}} \int_{x_{\min}}^{x_{\max}}
\frac{\mathrm{BA}(x) - c}{1 - c}\, dx,
\end{equation}
with $x = \ln(1 + k)$, chance level $c$ and the trapezoid rule over the
budgets of the test, so that chance gives 0 and a perfect decoder 1. We also report $k^*$, the
smallest budget at which a patient reaches 60 or 70\,\%. Paired comparisons
use the Wilcoxon signed-rank test, with the rank-biserial correlation $r$ and
95\,\% bootstrap confidence intervals (CIs) over patients (10\,000
resamples). As pre-registered, H1, H1b, H2b and H3c are one-sided and H2a and
H3 two-sided; $W$ is the sum of positive ranks for a one-sided test and the
smaller rank sum for a two-sided one. Holm's method corrects within each
pre-registered family.

The pre-registration fixed six hypotheses before the main study. Besides H1
and H1b, H2a compares the ALC over $k \in \{1, 2, 3, 5, 10\}$ of transfer and
from scratch for a new user on three classes. H2b states that transfer's
advantage is smaller with the whole pool than at $k = 5$. H3 makes the ALC
comparison, over $k \in \{1, 2, 3, 5\}$, for a new command on two co-primary
outcomes: (a) sixth finger against rest and (b) three classes. H3c states
that transfer detects the sixth finger better than the own-hand decoder at
$k = 2$ and $k = 5$. Each test ends in a pre-registered category:
\emph{supported} if its rule holds; \emph{contrary}, possible only for a
two-sided test, if the effect is significant in the other direction;
\emph{partial} for a two-part test that passes one part; and otherwise
\emph{inconclusive}, which is never read as evidence of no effect. A
one-sided test whose effect lies in the other direction is therefore
inconclusive. The stroke tests were repeated without the two development
patients and on the post-therapy session (also with sub-21 and sub-22
averaged into one value), and the calibration tests with a chronological
draw.

\subsection{Pre-registration, deviations and prior knowledge}
\label{sec:prereg}
The pre-registration was frozen in the project's git repository under the
tag \texttt{prereg-v1}, on a commit dated 3 October 2026 (the tag itself was
made on 8 October), before any decoder was trained on a patient other than
the two development patients, apart from one transfer smoke test. These are git's own dates: the plan was not
deposited with an independent registry. The frozen plan and its deviation
log are available from the authors on request; the supplement lists the
log's \nDeviations{} dated entries.
Most entries clarify a rule before any calibration result, such as how
$k = \text{all}$ draws. Before the confirmatory results the networks were
re-run, after we found that they had skipped the last incomplete batch of
each epoch. After them, the planned electrode-count study was dropped, a
sanity check whose premise proved wrong was recorded as failed, and the
analyses marked post hoc were added; the follow-ups among them were fixed in
settings files before they were run.

Before the plan was frozen we had analysed every patient and session with
labels (spectra, desynchronisation maps), and decoding on the development
patients had shown the direction of H1. Before the calibration study, the
whole-pool values of the from-scratch methods, which enter H2b, and the
pre-trained EEGNet's accuracy without calibration were known from the
benchmark. Exploratory tests are uncorrected, except that mirroring was
Holm-adjusted across budgets as planned; none changes a confirmatory
category.

The analysis uses MNE~\cite{gramfort2013}, pyRiemann~\cite{pyriemann},
scikit-learn~\cite{pedregosa2011}, braindecode~\cite{braindecode} and
PyTorch~\cite{paszke2019}; every result row records the code commit that
produced it. The continuous replays use the command gate of a demonstrator
that moves a virtual hand with a sixth finger.
